\scanpage{5}\mspage{2}
la structure de torseur sur cet ens, on considère le sous-groupe abélien distingué \deleted{\unsure{et}} non nul $V$ de $G$, de sorte que $V$ \deleted{\unsure{soit}} est un vectoriel de rang $2$ sur $\mathbb{F}_2$, et $G/V\xrightarrow{\sim}\mathrm{Aut}_{\mathbb{F}_2}(V)$, et on note que les \inserted{3-}sous-groupes de Sylow de $G$ correspondent aux ``relèvements'' de l'unique 3-sous-groupe de Sylow de $G/V\simeq\unsure{\mathfrak{S}|G}_J$. ($J=V^\times$), savoir $\mathfrak{S}_J^+=[\mathfrak{S}_J,\mathfrak{S}_J]$, d'indice $2$ dans $\mathfrak{S}_J$. \deleted{\unsure{Comme}} L'ens. $I_G$ des 3-sous-groupes de Sylow de $G$ s'identifie donc à l'ens des scindages de l'extension
\[
\cdot\quad\unsure{1\text{---}|\vdash|1\longrightarrow} V\unsure{\text{---}|\longrightarrow} G^+\longrightarrow\mathfrak{S}_J^+\unsure{\text{---}|\longrightarrow}1
\]
du 3-groupe $\mathfrak{S}_J^+$ par le 2-groupe $V$. Sur l'ens de ses scindages, $V$ opère donc transitivement par autom. intérieurs, d'ailleurs $V^{\mathfrak{S}_J^+}=\{0\}$ ($\mathfrak{S}_J^+$ est le groupe des permutations circulaires sur $V^\times=J$), donc $I_G$ est bien un torseur sous $V$.

Quand on part de $\unsure{J|I}$ et \deleted{$\unsure{\mathfrak{S}}$} $G=\mathfrak{S}_I$, on trouve une bijection
\[
\begin{aligned}
I&\unsure{\text{---}|\longrightarrow} I_G\\
i&\longmapsto\mathfrak{S}_{I\setminus\{i\}}^+=\text{l'unique 3-ss-gr. de Sylow}\\
&\hspace{8em}\text{de }\mathfrak{S}_I\text{ qui fixe }i
\end{aligned}
\]
Ainsi \deleted{\unsure{les deux foncteurs}} \inserted{les} foncteurs $I\unsure{\text{---}|\longrightarrow|\longmapsto}\mathfrak{S}_I$, $G\unsure{\text{---}|\longrightarrow|\longmapsto} I_G$, et la comparaison $I\longrightarrow I(\mathfrak{S}_I)$, on a un iso avec le \deleted{l'identité} foncteur identique,
% archive sheet number: 4
